\section{Transporting a local representation through one scale}
\label{sec:local-transfer}
We now explain the common transport argument used by the centered
estimate and the local OPE construction. Its analytic input is a
bound for integrating finitely many elementary insertions against
one fixed background transition. The argument below turns that
bound into a transition for arbitrary local fields represented by
absolutely summable elementary insertions. It also restores the
prescribed singleton profiles after integration. The same background
transition is used throughout.

\subsection{From an anchored estimate to a site estimate}
The first estimate accounts for summing over all named sources.
Write
\[
             S_M=e\{2(16M+9)\}^4.
\]
Suppose at most $N$ marks are anchored at any one site. Let
$X_i$ be supports, $A_i$ their nonempty sets of marks, and
$f_i\ge0$. Every $i$ with $f_i>0$ is assumed to have some
$v\in A_i$ with $o_v\in X_i$. If
\[
       \sum_{i:v\in A_i}e^{(a+1)c_M(X_i)}f_i\le c
                         \quad\text{for every }v,
\]
where $c\ge0$, then
\begin{equation}\label{eq:transfer-root-hit}
       \sum_{i:x\in X_i}e^{a c_M(X_i)}f_i\le NS_Mc
                         \quad\text{for every }x.
\end{equation}
Indeed, the defining connected-cover bound gives, for $x,y\in X$,
\[
       \operatorname{dist}(x,y)\le(4M+2)c_M(X)+2M.
\]
Hence a term on the left of \eqref{eq:transfer-root-hit} is at most
\[
 \sum_{v\in A_i}
 e^{1-\operatorname{dist}(x,o_v)/(4M+2)}
                         e^{(a+1)c_M(X_i)}f_i.
\]
Retain the anchored $v$ to verify this inequality; all other terms
are nonnegative. Sum first over $i$, use the anchored bound, and
then the occupancy bound. It remains to check
\[
 \sum_y e^{-\operatorname{dist}(x,y)/(4M+2)}
                  \le\{2(16M+9)\}^4.
\]
Lift $y-x$ to its minimal integer representative. These lifts are
distinct, and their $\ell^1$ norm is at most four times their
periodic maximum distance. With $b=1/(4(4M+2))$, the sum is
therefore bounded by
\[
 \sum_{z\in\mathbb Z^4}e^{-b|z|_1}
 =\left(1+\frac2{e^b-1}\right)^4
 \le\{2(1+b)/b\}^4=\{2(16M+9)\}^4.
\]
This proves the estimate, uniformly in the period. The argument
also proves summability for countable activity families by bounding
every finite subsum.

\subsection{Preparing a finite family of insertions}
Consider a fine representation $R$ and a coarse representation $E$.
Let $\mathcal T$ be their block-integration operator. For a fine
elementary insertion $M$, write $\mathcal Q(M)$ for its quoted
coarse flagged profile. We assume this map is real linear under
absolutely summable combinations in a fixed slot, with a uniform
bound by the elementary insertion cost. These are properties of
the fixed transition; their analytic construction is not being
asserted here.

For each mark $v$, suppose we have an anchored expansion, using
only finitely many distinct slots but possibly countably many
functions in each slot,
\begin{align*}
 g_v-u_v&=\sum_j z_{vj}\,\operatorname{ord}(M_{vj}),&
 q_v&=\sum_j z_{vj}\,\operatorname{flag}(M_{vj}),\\
 q'_v&=\sum_j z_{vj}\,\mathcal Q(M_{vj})
                                      +\mathcal H_E(\psi_v).
\end{align*}
Here $\sum_j|z_{vj}|\operatorname{cost}(M_{vj})\le y_v$,
and $\psi_v(X,U)$ is a coarse response anchored at the parent of
$o_v$. It has local dependence, is measurable and gauge invariant,
and vanishes when $X$ is healthy. Its cost hypothesis is
\begin{equation}\label{eq:transfer-response-budget}
 \sum_X e^{(p+1)c_{M_E}(X)}w_v\sup_U|\psi_v(X,U)|\le W.
\end{equation}
An available majorant can replace the displayed supremum. The
operator $\mathcal H_E$ will be described below. The constant
$u_v$ is outside the priced sum.

Compress the countable sum separately in each slot. Its real and
imaginary parts converge uniformly, since their absolute values
are bounded by the summable insertion weights. The sums are real
measurable local gauge-invariant functions in that same slot.
Their bounds are the respective sums of absolute coefficient
weights. Thus two real insertions, multiplied by $1$ and $\mathrm i$,
replace the whole complex series. Both the old profiles and their
quoted profiles are preserved, by the assumed continuity and
linearity. The resulting finite list for mark $v$ has weighted
cost at most $2w_vy_v$. If $2w_vy_v\le d_0$ and the fine-site
occupancy is at most $N_f$, \eqref{eq:transfer-root-hit} bounds its
site price at exponent $a$ by $N_fS_{M_R}d_0$, provided
$a+1\le\nu_R$.

Suppose the higher ordinary and flagged arrays in the fine local
representation each have site price at most $q$ at exponent $a$.
The real/imaginary construction of Section~\ref{sec:ope-singletons},
now applied to every coefficient of order at least two, represents
them by finite elementary insertions at total site price at most
$4q$. The unused-slot property is used here. Combining the two
lists gives a finite list with price
\begin{equation}\label{eq:transfer-inventory-cost}
                       T=N_fS_{M_R}d_0+4q.
\end{equation}
Its ordinary sum is $\sum_v(g_v-u_v)e_{\{v\}}+H$,
and its flagged sum is $\sum_vq_ve_{\{v\}}+K$, where $H,K$
are the higher arrays. All marks of each insertion lie in its support.

The elementary integration input takes a family of linear slot
insertions. To turn the exponential insertions in the last display
into that form, adjoin, at each ordinary slot, the product of every
compatible set of at least two members of the list. Compatibility
means that their mark sets are disjoint and their union is an
allowed monomial. Keep all original members. The added product
uses the same slot and support, and its price is the product of
the constituent prices. No product is added to a flagged slot.

This operation has price at most
\begin{equation}\label{eq:transfer-expanded-cost}
                              T+e^TT.
\end{equation}
To verify it, let $r_j\ge0$ be the product of the marker toll,
the modulus of the complex coefficient, and the elementary
supremum-norm weight, excluding the support exponential
$e^{a c_M(X)}$. For one ordinary slot put $s=\sum_jr_j$. Because every member has at least one
mark in that support and $a\ge0$, the site bound at any point of
the support implies $s\le T$. The sum of the product prices of
all nonempty subsets is at most
\[
       \prod_j(1+r_j)-1\le e^s-1\le se^s\le e^Ts.
\]
The added compatible subsets of size at least two form a subfamily.
Multiplying by the support exponential and summing over distinct
slots meeting a given site gives at most $e^TT$. Adding the
original price proves \eqref{eq:transfer-expanded-cost}.

The density identity is exact. For each ordinary slot, the sum
of all the compatible products is its squarefree exponential
minus one: expand $\prod_j(1+m_j)=\exp(\sum_jm_j)$, where each
nonempty monomial $m_j$ has square zero. Terms with incompatible
marks vanish in the algebra. Thus replacing the exponential at
that slot by one plus the expanded list changes no coefficient.
Flagged slots have not changed. Denote the resulting elementary
density by $\mathcal E$. The fine density $X$ satisfies
\begin{equation}\label{eq:transfer-elementary-identity}
             X=\exp\left(\sum_vu_ve_{\{v\}}\right)\mathcal E.
\end{equation}
Every added insertion has at least two marks, so the sum of quoted
singleton profiles is still $q'_v-\mathcal H_E(\psi_v)$.

\subsection{Absorbing the coarse response}
We state explicitly the nonlinear integration input used at this
point. For every prepared finite list of price at most $M_*$,
the operator $\mathcal T$ produces ordinary and flagged arrays
$\widehat F,\widehat K$, with the support and eligibility properties
of Section~\ref{sec:ope-singletons}, and
\begin{equation}\label{eq:transfer-nonlinear-input}
 \mathcal T\mathcal E
       =e^{E_E}\exp\left(\sum_X\widehat F_X\right)
                                       \mathcal G(\widehat K).
\end{equation}
The ordinary array has no constant term and site price at most
$k$ at exponent $t+\eta$; the flagged array has site price at
most $u$ at exponent $p$. Its vacuum coefficient is the coarse
background profile, and its singleton coefficient is precisely
the sum of the quoted singleton insertions. This estimate is
required uniformly in the finite list and period, subject to the
specified parent-site occupancy. Its production is a separate
analytic obligation.

Assume $T+e^TT\le M_*$. Require also that every prepared
elementary density is integrable under the fine prior and in
each transition slice. The lawful background representation and
its transition-kernel property supply these integrability
conditions in the application. Integration is componentwise on the
finite-dimensional polynomial algebra, so its continuous linear
maps commute with an integrable transition. In particular the
constant polynomial exponential in
\eqref{eq:transfer-elementary-identity} can be pulled out. We
obtain the same representation as
\eqref{eq:transfer-nonlinear-input}, with exponent
$P=\sum_vu_ve_{\{v\}}+\sum_X\widehat F_X$.

Set $J(U)=\sum_{v,X}\psi_v(X,U)e_{\{v\}}$. We next move $J$
from the ordinary exponent into the flagged array, without changing
the density. In a term of
\[
                   \mathcal G(\widehat K)\exp(J),
\]
regard each selected flagged pair and each selected response
insertion $(v,X)$ as a vertex, and join vertices whose supports
overlap. Every nonzero response insertion meets a flag, because it
vanishes on healthy supports. That flag is covered by a selected
flagged pair in the gas term. Hence every connected component
contains a flagged vertex. Within a component the flagged supports
are disjoint; give the component the union support and the union
of their flag sets. Define $\overline K$ as the sum of these
component products with that output pair, retaining its required
flag indicator. Every flag in a response support is covered by
an old flagged vertex meeting that support, hence by a vertex
in the same component. Thus the component flag set is exactly
the set of flags in its union support.

Connected-component decomposition is a bijection on nonzero
terms. In the forward direction, different components have disjoint
supports and their flag sets cover all flags. In the reverse
direction, the constituent vertices of components with disjoint
output supports cannot coincide or acquire a new overlap edge.
Their union recovers the original selection. The products of their
coefficients agree. Thus
\begin{equation}\label{eq:transfer-closure}
       \mathcal G(\overline K)=\mathcal G(\widehat K)\exp(J),
 \qquad e^P\mathcal G(\widehat K)
                           =e^{P-J}\mathcal G(\overline K).
\end{equation}
All these expansions are finite in support and mark indices.

With no response vertex a connected component is a single old
flagged vertex, since old flagged supports are disjoint. This
proves preservation of the vacuum. With exactly one response
vertex $(v,X)$, every flagged vertex in its component meets $X$:
there are no edges between distinct flagged vertices, so any path
to the response must be a direct edge. Let
$\Pi_E(X;Y,Z;U)$ be the sum of products of the background flagged
activities over these nonempty, pairwise-disjoint collections
meeting $X$, with union support $Y$ after adjoining $X$ and union
flag set $Z$, including the output flag indicator. Then
\[
 \mathcal H_E(\psi_v)(Y,Z;U)
                  =\sum_X\psi_v(X,U)\Pi_E(X;Y,Z;U).
\]
This defines the operator in the transition expansion above.
Two or more response vertices have degree at least two and cannot
contribute a singleton. Therefore
\[
 [e_{\{v\}}]\overline K
       =[e_{\{v\}}]\widehat K+\mathcal H_E(\psi_v)=q'_v.
\]
Replacing the background profiles in $\Pi_E$ by their countable
indexed sums is justified by absolute convergence, bounded by
the product of $1+|k_j^0|$, as in
\eqref{eq:ope-single-coin}.

The response insertions have site price at most
$N_cS_{M_E}W$ at exponent $p$, by
\eqref{eq:transfer-response-budget} and
\eqref{eq:transfer-root-hit}, where $N_c$ is the parent-site
occupancy. Put
\[
                              A=u+N_cS_{M_E}W.
\]
If
\begin{equation}\label{eq:transfer-closure-budget}
 t+\eta\ge0,\qquad
               (2M_E+1)^4A\le p-(t+\eta),
\end{equation}
the connected-support estimate \eqref{eq:ope-forest-price} bounds
the site price of $\overline K$ at exponent $t+\eta$ by $e^A-1$.
Its coefficients are products of old coefficient majorants, and
the map from a constituent set to its output pair is unique;
thus the counting lemma applies without a multiplicity factor.
The ordinary higher coefficients retain price $k$, since $J$ is
linear. Local dependence, gauge invariance, measurability and the
flag condition follow from finite products and support-local flag
indicators.

Finally suppose the singleton coefficient of the actual integrated
density has already been identified as the linear density of a
prescribed pair $(g'_v,k'_v)$, with $k'_v$ dominated and profile
$q'_v$. Assume the coarse unmarked density is nonzero. The
uniqueness following \eqref{eq:ope-single-coin} identifies the
ordinary singleton of $P-J$ with $g'_v$. Removing all coefficients
of degree zero and one now yields a local representation with
these prescribed fields, with ordinary price $k$ and flagged price
$e^A-1$ at exponent $t+\eta$.

\subsection{Changing the mark weights}
The final operation pays for changing weights. Suppose $w_v,w'_v>0$
and for every allowed monomial $B$ with $|B|\ge2$ supported in $X$,
\begin{equation}\label{eq:transfer-reweight}
       \prod_{v\in B}w'_v
           \le a_*\prod_{v\in B}w_v\,e^{\eta c_{M_E}(X)},
                         \qquad a_*\ge0.
\end{equation}
Multiply this inequality by each coefficient majorant and sum.
It changes an exponent-$t+\eta$ price $C$ into an exponent-$t$
price at most $a_*C$. Hence the transported representation has
both higher-array prices bounded by $g$ whenever
\[
                       a_*k\le g,\qquad a_*(e^A-1)\le g.
\]

Here is the specialization used in the scale induction. Start
with weights $I_v>0$, higher-array price $1$, and exponent at least
$\nu_R-4$. Choose $0<d\le1$ and first reduce the weights to
$w_v=dI_v$. A monomial with at least two marks acquires the factor
$d^{|B|}\le d$, so the higher-array price becomes $q=d$.
Suppose the singleton transition costs obey
$2I_vy_v\le1$ and $I_vr_v\le1$, with response cost $r_v$ measured
at the full coarse exponent plus two. After multiplying by $w_v$,
the fine singleton budget is $d_0=d$ and the response budget is
$W=d$ at $p=\nu_E+1$.

Assume $\nu_R,\nu_E\ge100$, and write $N$ for a bound on both
occupancies. The two numerical conditions are
\begin{align}
 T&=NS_{M_R}d+4d,& T+e^TT&\le M_*,\label{eq:transfer-duty}\\
 A&=u+NS_{M_E}d,& (2M_E+1)^4A&\le\tfrac12,
                         &e^A-1&\le1.\notag
\end{align}
Use the nonlinear input with fine exponent $\nu_R-4$, ordinary
output exponent $\nu_E+1/2$ and price $1$, and flagged output
exponent $\nu_E+1$ and price $u$. The fine anchored estimate
requires only $\nu_R-3\le\nu_R$, and all exponents used for
support counting are nonnegative. Equation~\eqref{eq:transfer-duty}
therefore verifies \eqref{eq:transfer-closure-budget} with
$t=\nu_E$ and $\eta=1/2$. Any new weights satisfying
\eqref{eq:transfer-reweight} with $a_*=1$ give a coarse local
representation at exponent $\nu_E$ and price $1$.
This proves the one-scale propagation from the stated elementary
transition estimate and singleton transition expansions. It does
not replace the proof of that analytic estimate or the construction
of those expansions.
